% =========================================================================== %
%  P2 (BASKET) pad-Micromegas cosmic-bench characterisation report
%  Detectors P2_1 and P2_2  —  CEA Saclay cosmic test bench, June–July 2026
%
%  Build:  latexmk -pdf p2_cosmic_bench_report.tex
%  Figures are included directly from the Analysis/ tree next to the data
%  (see \graphicspath below); re-running the pipeline refreshes them.
% =========================================================================== %
\documentclass[11pt,a4paper]{article}

\usepackage[margin=2.4cm]{geometry}
\usepackage[T1]{fontenc}
\usepackage[utf8]{inputenc}
\usepackage{lmodern}
\usepackage{amsmath}
\usepackage{graphicx}
\usepackage{booktabs}
\usepackage{siunitx}
\usepackage{xcolor}
\usepackage{caption}
\usepackage{subcaption}
\usepackage{enumitem}
\usepackage[hidelinks]{hyperref}
\usepackage{microtype}

\graphicspath{{/local/home/ak271430/Documents/PostDocSaclay/data/Cosmic_Bench/Analysis/}}

\captionsetup{font=small,labelfont=bf}
\sisetup{separate-uncertainty=true}

\newcommand{\detone}{P2\_1}
\newcommand{\dettwo}{P2\_2}
\newcommand{\code}[1]{\texttt{#1}}

\title{\textbf{Characterisation of the P2 (BASKET) pad-Micromegas detectors\\
       on the CEA Saclay cosmic test bench}\\[1ex]
       \large Detectors \detone{} and \dettwo{} --- runs of 30 June -- 10 July 2026}
\author{Alexandra Kallitsopoulou\\ \small CEA Saclay, IRFU/DEDIP}
\date{11 July 2026}

\begin{document}
\maketitle

\begin{abstract}
\noindent
We report the first cosmic-ray characterisation of two P2 (BASKET) detectors ---
metallic, non-resistive pad-readout Micromegas --- on the CEA Saclay cosmic test
bench, using the four-plane M3 Micromegas telescope as an external tracking
reference. A fully automated DAQ and processing chain (DREAM front-end readout,
on-the-fly decoding and hit reconstruction, M3 track reconstruction, HV
monitoring) feeds a staged offline QA pipeline that validates the
Gerber-derived pad mapping, aligns each detector to the telescope, vetoes
HV discharges, and measures track-referenced efficiency and position response.
Detector~2 reaches a validated position-reconstruction efficiency of
\SI{65}{\percent} integrated over a \SI{10}{h} run (\SIrange{74}{75}{\percent}
in its best HV points), with a core position resolution of
$\sigma \simeq \SI{4.7}{mm}$, limited by the \(\sim\SI{12}{mm}\) pad pitch.
Detector~1 showed an efficiency of \SI{56}{\percent} in its first long run but
suffers from severe intermittent gain drops whose duty cycle degraded across
subsequent runs; understanding and curing this instability is the top priority
of the next campaign. The efficiency plateau has not yet been reached in either
mesh-HV scan. We describe the measurement, the trigger and DAQ logic, the
processing pipeline and the analysis methodology in detail, and conclude with a
roadmap for the coming measurements.
\end{abstract}

\tableofcontents
\clearpage

% ========================================================================== %
\section{Introduction}
% ========================================================================== %

The P2 ``BASKET'' detector is a \emph{metallic, non-resistive} bulk Micromegas
with a segmented \emph{pad} anode readout. Unlike the resistive-strip
Micromegas studied previously on the same bench, the P2 anode is divided into
1280 pads arranged in a fan-shaped (basket) sector: 10 physical connector
sectors of 128 pads each, with a typical pad size of order
$12\times12\,\text{mm}^2$. Each pad carries a full two-dimensional position, so the
spatial analysis unit is the pad itself --- there is no strip clustering or
micro-TPC reconstruction. The pad geometry (centre, polar coordinates, area) is
derived directly from the production Gerber files of the readout PCB.

Two detectors of this design, \detone{} and \dettwo{}, are currently installed
on the cosmic-ray test bench at CEA Saclay. The goals of this first
characterisation campaign are:

\begin{enumerate}[itemsep=2pt]
  \item validate the full channel-to-pad mapping chain (Gerber map $+$ DREAM
        front-end wiring) against reconstructed cosmic tracks;
  \item measure the track-referenced detection efficiency and its uniformity
        across the pad plane, including the effect of the mesh-support
        insulation pillars;
  \item measure the position response (residuals, core resolution) of the pad
        centroid reconstruction;
  \item characterise the discharge (spark) behaviour of the metallic mesh and
        establish clean veto procedures;
  \item map the efficiency and resolution as a function of the mesh
        amplification voltage (HV scan) to locate the operating plateau.
\end{enumerate}

This note documents the measurement setup
(Section~\ref{sec:setup}), the trigger and DAQ logic
(Section~\ref{sec:trigger}), the automated processing pipeline
(Section~\ref{sec:pipeline}), the offline analysis methodology
(Section~\ref{sec:analysis}), the results for both detectors
(Section~\ref{sec:results}), and the conclusions and roadmap
(Sections~\ref{sec:conclusions} and~\ref{sec:roadmap}).

% ========================================================================== %
\section{Measurement setup}
\label{sec:setup}
% ========================================================================== %

\subsection{Bench geometry}

The cosmic bench is a vertical tower. Cosmic muons are tracked by the
\emph{M3 reference telescope}: four double-coordinate (X/Y strip) resistive
Micromegas planes, two below and two above the detectors under test.
The two P2 detectors sit between the telescope stations. The nominal plane
heights from the run configuration are listed in
Table~\ref{tab:geometry}. For \dettwo{} the plane height fitted from the data
($z = \SI{712}{mm}$, Section~\ref{sec:alignment}) is used in the analysis
instead of the nominal \SI{702}{mm}.

\begin{table}[htb]
  \centering
  \caption{Bench elements and their vertical positions ($z$, measured upwards
  along the tower). HV channels refer to CAEN crate \code{card:channel}.}
  \label{tab:geometry}
  \small
  \begin{tabular}{llrl}
    \toprule
    Element & Type & $z$ [mm] & HV (mesh / drift) \\
    \midrule
    \code{m3\_bot\_bot} & M3 X/Y plane & 24   & 3:8 / 0:8 \\
    \code{m3\_bot\_top} & M3 X/Y plane & 144  & 3:9 / 0:9 \\
    \detone{}           & P2 pad Micromegas & 232 & 1:0 / 1:1 \\
    \dettwo{}           & P2 pad Micromegas & 702 (fitted 712) & 1:2 / 1:3 \\
    \code{m3\_top\_bot} & M3 X/Y plane & 1185 & 3:10 / 0:10 \\
    \code{m3\_top\_top} & M3 X/Y plane & 1302 & 3:11 / 0:11 \\
    \bottomrule
  \end{tabular}
\end{table}

\subsection{Detectors under test}

Each P2 detector is read out through 10 physical 128-pin connectors
(``sectors''), each split into a \emph{bottom} half (pads 1--64) and a
\emph{top} half (pads 65--128); every half is wired to one 64-channel DREAM
front-end connector. The full detector therefore maps to 20 DREAM connectors
spread over three front-end units (FEUs). The gas is
Ar/iC$_4$H$_{10}$ 95/5 at atmospheric pressure. Two HV channels per detector
power the amplification mesh and the drift electrode; the drift gap is
operated at a fixed $V_\text{drift} - V_\text{mesh} = \SI{170}{V}$ during the
combined runs.

Not all connectors were live during the campaign, and the analysis explicitly
removes dead connectors from both the numerator and the denominator of every
efficiency estimate:

\begin{itemize}[itemsep=2pt]
  \item \detone{}: connector~10 disconnected in the first run; connectors 1
        and 10 disconnected from run 2 onwards (9 or 8 live sectors,
        1152/1024 pads);
  \item \dettwo{}: connector~1 physically disconnected; connectors 8--10 are
        wired to FEU~8, which crashes the DAQ and is excluded from the runs.
        Live readout is connectors 2--7 (768 pads), confirmed independently
        by track--hit correlation.
\end{itemize}

The insulation mask Gerber of the mesh-support pillars
($\sim$11\,700 small $\sim$\SI{0.5}{mm} pillars on a \SI{4}{mm} grid plus five
large $\sim$\SI{3.3}{mm} ones) is loaded in the same PCB frame as the pad map
and overlaid on all hit and efficiency maps, so that pillar shadows can be told
apart from genuine gain defects.

\subsection{Runs analysed}

\begin{table}[htb]
  \centering
  \caption{Long runs and HV scans entering this report. All runs use
  Ar/iC$_4$H$_{10}$ 95/5. Mesh/drift voltages in volts.}
  \label{tab:runs}
  \footnotesize
  \setlength{\tabcolsep}{3pt}
  \begin{tabular}{lcccl}
    \toprule
    Run & Detector(s) & Mesh/drift & Duration & Purpose \\
    \midrule
    \code{p2\_det1\_long\_run\_6-30-26} & \detone{} & 440/600 & long & first efficiency run \\
    \code{p2\_det1\_mesh\_hv\_scan\_7-2-26} & \detone{} & 345--420, 5\,V steps & scan & HV scan \\
    \code{p2\_det1\_long\_run\_7-4-26} & \detone{} & 440/600 & long & repeat \\
    \code{p2\_det1\_long\_run\_7-7-26} & \detone{} & 440/600 & long & repeat \\
    \code{p2\_det1\_det2\_long\_run\_mesh\_scan\_7-9-26} & \detone{}$+$\dettwo{} & 430/600 & 10 h $+$ scan & first \dettwo{} run \\
    \bottomrule
  \end{tabular}
\end{table}

The 9 July combined run is the first with \dettwo{} installed: a \SI{10}{h}
long sub-run at mesh \SI{430}{V} followed by a mesh scan from 430 down to
\SI{395}{V} in \SI{5}{V} steps, with the drift stepped in tandem to keep the
drift gap voltage constant at \SI{170}{V}.

% ========================================================================== %
\section{Trigger logic and data acquisition}
\label{sec:trigger}
% ========================================================================== %

\subsection{Trigger}

The bench runs on an \emph{external cosmic trigger}: the DREAM system is
configured with \code{Sys DaqRun Trig = Ext} and trigger source
\code{Tg\_Src\_ExtSyn}, i.e.\ a common trigger signal is distributed
synchronously to all front-end units through the trigger interface. The
trigger FEU is the M3 telescope FEU (FEU~1, set from \code{m3\_feu\_num} in
the run configuration), so the trigger definition is tied to the reference
telescope rather than to the detectors under test; top and bottom trigger
scintillator counters are installed on the tower (above the top and below the
bottom M3 stations) for trigger formation. Two consequences shape the
analysis:

\begin{itemize}[itemsep=2pt]
  \item The trigger is \emph{unbiased} with respect to the P2 detectors: a
        triggered muon crossing a P2 plane that produces no P2 hit is a
        genuine inefficiency and is kept in the efficiency denominator.
  \item P2 discharges are only rarely recorded by the DAQ (they are not
        time-correlated with cosmic triggers), so the spark treatment must be
        driven by the HV monitoring rather than by the event data alone
        (Section~\ref{sec:sparks}).
\end{itemize}

On every accepted trigger, all active FEUs read out 32 ADC samples per channel
(no zero suppression in these runs); every hit carries the global trigger
timestamp (\code{trigger\_timestamp\_ns}) and event number, which are used
offline to match P2 hits, M3 tracks and HV records in time.

\subsection{DAQ orchestration}

Data taking is orchestrated by the \code{Cosmic\_Bench\_DAQ\_Control} package,
a client--server system over TCP sockets with persistent \code{tmux}-managed
servers:

\begin{itemize}[itemsep=2pt]
  \item \textbf{HV server} (\code{hv\_control.py}): ramps and monitors the
        CAEN crate; during runs it logs $v_\text{mon}$/$i_\text{mon}$ of every
        channel to \code{hv\_monitor.csv} at sub-second cadence
        ($\sim$\SI{0.5}{s} interval, $\sim$\SI{2}{s} effective granularity).
  \item \textbf{DREAM DAQ server} (\code{dream\_daq\_control.py}): writes the
        run-specific DREAM configuration (active FEUs and connectors derived
        from the declared detectors, trigger FEU preserved), launches
        pedestal-threshold and data runs, and copies raw FDF files on the fly.
  \item \textbf{M3 tracking server} (\code{m3\_tracking\_control.py}): runs
        the M3 telescope track reconstruction as files arrive.
  \item \textbf{Orchestrator} (\code{daq\_control.py}): sequences sub-runs ---
        for each sub-run it sets the HV working point, waits for ramp
        completion, starts the DREAM run for the configured duration, and
        stops it; at the end of the session HV is powered down.
  \item \textbf{Processor and QA watchers}: autonomous processes that decode,
        reconstruct and QA-plot the data as it arrives
        (Section~\ref{sec:pipeline}).
\end{itemize}

Each run is fully described by a \code{run\_config.json} snapshot (detector
geometries, FEU wiring, HV set points per sub-run, DAQ settings), which the
offline analysis treats as the single source of truth for wiring and
geometry.

Before every data run a dedicated \emph{pedestal-threshold run} measures the
per-channel pedestal mean and RMS. The DREAM waveform processing then keeps a
pulse only if its amplitude exceeds $5\,\sigma_\text{ped}$ of its own channel,
so each channel's effective threshold is set by its own pedestal noise --- a
fact exploited by the pedestal QA stage (Section~\ref{sec:pedestal}).

% ========================================================================== %
\section{Data processing pipeline}
\label{sec:pipeline}
% ========================================================================== %

\subsection{Online (on-the-fly) processing}

The \code{processor\_watcher.py} daemon drives a fully automatic chain as raw
files land on disk:

\begin{enumerate}[itemsep=2pt]
  \item \textbf{Decode}: raw DREAM FDF files $\rightarrow$ ROOT waveform trees
        (\code{decoded\_root/});
  \item \textbf{Waveform analysis}: per-channel pulse finding with the
        $5\sigma$ pedestal threshold $\rightarrow$ per-FEU hit trees
        (\code{hits\_root/}), carrying amplitude, time, saturation flag and
        the applied pedestals;
  \item \textbf{FEU combination}: per-FEU hits merged by event
        $\rightarrow$ \code{combined\_hits\_root/}, the main analysis input;
  \item \textbf{M3 tracking}: the M3 FEU data are reconstructed by the
        dedicated tracking code (v2) into per-event rays
        $\rightarrow$ \code{m3\_tracking\_root/} (positions at top/bottom
        stations, per-coordinate cluster counts and $\chi^2$);
  \item \textbf{HV log}: \code{hv\_monitor.csv} accumulates in parallel.
\end{enumerate}

A QA watcher produces online per-detector occupancy and rate plots for live
monitoring; a Flask dashboard exposes run status, HV and QA plots.

\subsection{Offline QA pipeline}

The offline analysis lives in \code{P2\_basket\_analysis/cosmic\_bench\_analysis}
and is organised as numbered, idempotent stages driven by a central run
registry (\code{p2\_qa\_config.py}); each stage writes its products into an
\code{Analysis/$\langle$det$\rangle$/$\langle$run$\rangle$/$\langle$sub\_run$\rangle$/$\langle$stage$\rangle$/}
tree that lives next to the data, organised per detector. The full pipeline
for a long run is reproduced by \code{run\_p2\_pipeline.sh}; every
veto-capable stage runs twice (with and without the discharge veto) and all
products carry explicit filename suffixes recording the dropped dead
connectors and the veto state, so no variant can silently overwrite another.
The stages are:

\begin{description}[itemsep=2pt,leftmargin=2.6cm]
  \item[02 mapping] pad-map validation against data occupancy, including a
        side-by-side comparison of candidate channel orderings;
  \item[03 alignment] P2$\leftrightarrow$M3 rotation scan, Procrustes
        transform, detector-plane $z$ scan;
  \item[04 M3 QA] reference-telescope quality (track $\chi^2$, multiplicity,
        angles, illumination), independent of the P2;
  \item[05 deep QA] hitmaps, per-pad firing fractions, event multiplicity and
        its time evolution, hot-pad hunting;
  \item[06 efficiency] track-referenced efficiency maps, breakdown and
        residuals; per-ray hit/miss list; efficiency vs.\ time;
  \item[07 HV sparks] mesh-discharge detection from the HV log, veto
        validation, DAQ cross-check;
  \item[08 pad sparks] per-pad micro-discharge flagging (review only, not
        masked);
  \item[09 pedestals] applied pedestal mean/RMS maps, threshold--rate
        correlation;
  \item[10 sliding map] kernel-smoothed (fixed and adaptive k-NN) efficiency
        maps at \SI{5}{mm} scale, resolving the large insulation pillars;
  \item[11 HV scan] efficiency and resolution vs.\ mesh HV across scan
        sub-runs, with a single frozen geometry;
  \item[12 validation] structural bias tests of the efficiency chain
        (Section~\ref{sec:validation}).
\end{description}

Two report builders (\code{build\_final\_pdf.py},
\code{build\_hv\_scan\_pdf.py}) compile all stage products into per-run
narrative PDF documents.

% ========================================================================== %
\section{Analysis description}
\label{sec:analysis}
% ========================================================================== %

This section describes the analysis at the algorithm level. Throughout, an
\emph{event} is one external trigger: the M3 telescope hits, the P2 hits and
the HV record are associated purely by the event number (\code{eventId}) and
the global trigger timestamp (\code{trigger\_timestamp\_ns}), which counts
nanoseconds continuously from run start across all combined-hits files. The HV
monitor also starts at run start, so the two clocks share a common $t=0$; the
agreement of the two time spans was verified to $\sim$\SI{0.3}{\percent}.

\subsection{Inputs and event model}

Each analysis stage consumes up to four inputs per sub-run:

\begin{itemize}[itemsep=2pt]
  \item \code{combined\_hits\_root/}: one row per fired channel per event,
        carrying the event number and trigger timestamp, the FEU and channel
        address, the pulse amplitude (ADC) and time (sample), and a
        saturation flag. Only channels above their individual
        $5\sigma_\text{ped}$ threshold appear, and \emph{only events with at
        least one hit anywhere in the DAQ are written} --- an important
        detail for the efficiency denominator (Section~\ref{sec:efficiency});
  \item \code{m3\_tracking\_root/}: per-event reconstructed rays with line
        endpoints $(X,Y)$ at two reference heights $Z_\text{Up}$,
        $Z_\text{Down}$, per-coordinate $\chi^2$ and cluster counts;
  \item \code{hv\_monitor.csv}: wall-clock time series of $v_\text{mon}$,
        $i_\text{mon}$ and power state for every CAEN channel
        ($\sim$\SI{2}{s} sampling);
  \item \code{hits\_root/} (per-FEU, pre-combination): additionally carries
        the \code{pedestals} tree (per-channel mean and RMS actually applied
        by the processing), used only by the pedestal QA.
\end{itemize}

\subsection{Channel-to-pad mapping and its validation}
\label{sec:mapping}

A DAQ hit is addressed by (FEU, channel), $\text{channel} \in [0,255]$ on a
4-connector DREAM front-end. The decode chain to a physical pad is:
%
\begin{align}
  \text{dream\_conn} &= \lfloor \text{channel}/64 \rfloor + 1, &
  \text{within} &= \text{channel} \bmod 64, \nonumber\\
  \text{(feu, dream\_conn)} &\;\xrightarrow{\;\text{run\_config wiring}\;}\;
  \code{c\_$N$\_bot/top}, &
  \text{sector} &= N-1, \nonumber\\
  \text{strip} &= \text{half\_base} + f(\text{within}), &
  \text{channel\_id} &= 128\,\text{sector} + (\text{strip}-1),
\end{align}
%
where \code{half\_base} is 1 for the \code{bot} half (strips 1--64) and 65
for the \code{top} half (strips 65--128), and the final
\code{channel\_id} keys the Gerber-derived map row carrying the pad centre
$(x,y)$, polar coordinates, tile polygon and area. The wiring dictionary is
taken from the run's own \code{run\_config.json} and inverted, so a rewiring
between runs is picked up automatically; pads on dead connectors are dropped
from the table at this point.

The only genuinely uncertain link is the within-half ordering function $f$:
whether DREAM channel 0 of a half corresponds to strip 1 or strip 64, or some
permutation. It is kept as an explicit, swappable strategy:
$f_\text{linear}(w) = w$, $f_\text{reverse}(w) = 63-w$, and
$f_\text{pairswap}$ (linear with adjacent pairs swapped, the pattern seen in
the MEC8 pin table). Flat cosmic illumination cannot distinguish these ---
all three fill the fan uniformly --- so the discrimination comes from the
\emph{per-event} correlation between the P2 pad centroid and the projected M3
track: with \code{reverse} the correlations reach $r_x = 0.93$,
$r_y = 0.88$ with a clean $\sim$\SI{89}{\degree} rotation and scale
$\simeq 1$, while \code{linear} and \code{pairswap} collapse to
$r_y \sim 0.15$. Since the within-half ordering is a fixed property of the
DREAM readout, \code{reverse} is the validated default for all P2 detectors
(re-confirmed on \dettwo{}).

Stage 02 renders the occupancy, mean amplitude and dead-pad maps on the
physical pad plane and checks three qualitative signatures of a correct map:
contiguous illumination (no speckle), dead connectors appearing as coherent
angular blocks, and radial/azimuthal profiles consistent with a uniform
cosmic flux over the fan (Fig.~\ref{fig:occupancy}).

\begin{figure}[htb]
  \centering
  \includegraphics[width=.92\textwidth]{det2/p2_det1_det2_long_run_mesh_scan_7-9-26/long_run_mesh_430V_drift_600V/02_map_validation/pad_occupancy_without_connectors_1_8_9_10_spark_vetoed.png}
  \caption{\dettwo{}, 9 July long run: cosmic hit occupancy projected on the
  physical pad plane (with radial/azimuthal profiles and per-connector
  occupancy). The illumination is contiguous over the live connectors 2--7;
  the missing sectors are the known dead connectors.}
  \label{fig:occupancy}
\end{figure}

\subsection{M3 reference tracks}
\label{sec:m3}

The M3 reconstruction (tracking v2) fits straight lines independently in the
$xz$ and $yz$ projections through the four telescope planes and stores each
ray as its endpoints at $Z_\text{Up}$ and $Z_\text{Down}$. The analysis
re-projects a ray to any plane $z$ by linear interpolation,
%
\begin{equation}
  x(z) = X_\text{Down} + \bigl(X_\text{Up}-X_\text{Down}\bigr)\,
         \frac{z - Z_\text{Down}}{Z_\text{Up} - Z_\text{Down}},
\end{equation}
%
and likewise for $y(z)$ --- which makes the alignment $z$-scan
(Section~\ref{sec:alignment}) essentially free.

A ray enters any P2 analysis only if it passes the quality recipe
%
\begin{equation}
  \chi^2_X,\;\chi^2_Y < 5
  \quad\text{and}\quad
  N^X_\text{clus},\;N^Y_\text{clus} \ge 3 ,
  \label{eq:m3cut}
\end{equation}
%
plus a single-track requirement (exactly one good ray in the event). The
cluster-count condition is not cosmetic: a fit with only two points per
coordinate is exact, so it carries a denormal-tiny (but non-zero) $\chi^2$
and would sail through any pure $\chi^2$ cut; requiring $\ge 3$ clusters
forces a genuine 3--4 plane fit. The $\chi^2$ threshold itself was re-tuned
for the v2 reconstruction (the old $\chi^2<20$ recipe belongs to the previous
$\chi^2$ scale and is stale). Stage 04 monitors the reference itself ---
$\chi^2$ distributions, good-track multiplicity, angular distributions
($\theta_x$, $\theta_y$ of clean single tracks), and the illumination
projected at the P2 plane --- so that a degraded telescope cannot silently
masquerade as a degraded P2.

\subsection{Per-event P2 observables}
\label{sec:centroid}

Hits are mapped to pads (Section~\ref{sec:mapping}); unmapped hits (dead
connectors, non-P2 FEUs) are discarded. For every event with $\ge 1$ mapped
pad the analysis computes the \emph{charge-weighted pad centroid}
%
\begin{equation}
  \mathbf{r}_\text{P2}
  = \frac{\sum_i a_i\, \mathbf{r}^\text{pad}_i}{\sum_i a_i},
  \qquad a_i = \max(A_i, 0),
  \label{eq:centroid}
\end{equation}
%
where the sum runs over the fired pads of the event, $A_i$ is the pulse
amplitude and $\mathbf{r}^\text{pad}_i$ the pad centre. The pad multiplicity
$n_\text{pad}$ is retained; genuine muon clusters have
$n_\text{pad}=1\text{--}3$ (median 1), which is what makes the
$n_\text{pad}\ge 20$ burst veto (Section~\ref{sec:sparks}) safe. A
\emph{leading-pad} estimator (position of the highest-amplitude pad) is
implemented as an alternative and used by the validation stage to bound the
cost of the centroiding choice (Section~\ref{sec:validation}).

\subsection{Alignment to the M3 reference}
\label{sec:alignment}

The M3 rays live in the bench frame (origin near the tower axis, mm); the pad
positions live in the PCB frame (a fan sector in the first quadrant of the
board coordinates). The transform between them --- translation
$\mathbf{t}$, rotation $\theta$ about $z$, possible reflection, and a scale
$s$ kept free as a diagnostic --- is determined from the data in three steps
(stage 03), on the sample of events that have both a clean single ray and a
P2 centroid.

\paragraph{1. Rotation scan.}
Pearson correlation coefficients are invariant under translation and scale,
so the rotation can be found \emph{without} first solving the translation.
The pad centroids are centred on their own mean and rotated by a trial angle
$\theta$,
\begin{equation}
  x'(\theta) = x_\text{pad}\cos\theta - \mu\,y_\text{pad}\sin\theta,\qquad
  y'(\theta) = x_\text{pad}\sin\theta + \mu\,y_\text{pad}\cos\theta,
\end{equation}
with mirror state $\mu = \pm 1$, and the figure of merit
\begin{equation}
  \mathrm{FOM}(\theta,\mu) = r\bigl(x_\text{M3}, x'\bigr)
                           + r\bigl(y_\text{M3}, y'\bigr)
\end{equation}
is scanned over $\theta \in [0,360)$\si{\degree} in \SI{0.5}{\degree} steps
for both mirror states (a reflection is not reachable within $SO(2)$, so it
must be tried explicitly). A correct mapping produces a single sharp peak
with both $r$ large and positive; a wrong channel ordering destroys the peak
in one coordinate, which is exactly the mapping test of
Section~\ref{sec:mapping}. Wild rays are suppressed by an M3 fiducial cut
($|x|,|y| < \SI{1500}{mm}$) and, optionally, a minimum pad amplitude.

\paragraph{2. Procrustes fit.}
Independently, the best-fit similarity transform is obtained in closed form
by SVD. With $P$ and $Q$ the mean-centred $N\times2$ matrices of pad and M3
positions,
\begin{equation}
  H = P^{\mathsf T} Q = U \Sigma V^{\mathsf T}, \qquad
  R = V U^{\mathsf T}, \qquad
  s = \frac{\operatorname{tr}\Sigma}{\lVert P \rVert_F^2}, \qquad
  \mathbf{r}_\text{M3} = s\,R\,(\mathbf{r}_\text{pad} - \bar{\mathbf{r}}_\text{pad})
                        + \bar{\mathbf{r}}_\text{M3},
  \label{eq:procrustes}
\end{equation}
with $\det R < 0$ signalling a reflection. The implied rotation angle must
coincide with the scan peak (it does, for both detectors), and the fitted
scale is a powerful health check: pads are physical objects, so a correct
mapping and geometry \emph{must} give $s \simeq 1$. The fit RMSE is dominated
by the pad pitch and by any gain intermittency mixing in poorly measured
events; only events inside a $\pm\SI{300}{mm}$ fiducial feed the fit.

\paragraph{3. Plane-height ($z$) scan.}
Because the M3 rays have slopes, the projected impact points depend on the
assumed plane height; the rotation/correlation scan is therefore repeated
while re-projecting the rays over a $\pm\SI{250}{mm}$ window around the
nominal $z$ (using the stored endpoints, Eq.~(2)). The height maximising the
correlation --- equivalently, the plateau of efficiency vs.\ $z$ in stage 12
--- defines the measured detector plane. For \dettwo{} this gives
$z = \SI{712}{mm}$ against a nominal \SI{702}{mm}; the fitted value is
registered in the run configuration (\code{PLANE\_Z}) and used by every
stage, with the rule that the measured plane always wins over the nominal
mounting height.

The shipped transform used by all downstream stages is the Procrustes fit,
re-fitted by the efficiency stage itself (Section~\ref{sec:efficiency}) on
its own matched sample --- stages share the \emph{same} fitting code
(\code{p2\_align.fit\_transform}), not frozen numbers, so no stage can drift
out of sync with the alignment.

\subsection{Discharge (spark) treatment}
\label{sec:sparks}

Two independent vetoes target two distinct discharge populations.

\paragraph{HV (whole-mesh) spark veto.}
The mesh baseline current is $\sim$\SI{0.07}{\micro A} (90th percentile
$\sim$\SI{0.14}{\micro A}); a discharge drives $i_\text{mon}$ to several
\si{\micro A}, up to the current-limit compliance ($\sim$\SI{10}{\micro A}),
while the crate keeps the channel powered. The veto is built from the HV log
as follows:

\begin{enumerate}[itemsep=2pt]
  \item flag every HV sample with $i_\text{mon} \ge \SI{2}{\micro A}$
        (well above baseline, well below compliance);
  \item group consecutive flagged samples into spark intervals, merging
        groups separated by $\le \SI{6}{s}$ (up to two missed \SI{2}{s}
        samples --- one flickering discharge, not two);
  \item integrate the released charge per spark,
        $Q = \int i_\text{mon}\,dt$ (trapezoidal, \si{\micro A.s} =
        \si{\micro C}), and record duration, peak current and sample count;
  \item pad each interval by \SI{-2}{s} (clock skew) and \SI{+10}{s}
        (post-discharge recovery, during which the gain is degraded), and
        merge overlaps;
  \item drop every event whose trigger time
        $t = \code{trigger\_timestamp\_ns}/10^9$ falls inside a padded
        interval.
\end{enumerate}

The vetoed live time is bookkept and reported
(\code{live\_fraction}); crucially, the veto is applied as a set of
\emph{event IDs} removed identically from the M3 ray list, the P2 centroid
list and the fired-event set, so the efficiency numerator and denominator
always see the same live sample. Stage 07 validates the construction: the
$v_\text{mon}$/$i_\text{mon}$ timeline with shaded veto windows, the spark
rate evolution, the per-spark charge spectrum, and a cross-check against DAQ
high-multiplicity bursts. Because the trigger is external, discharges are
rarely captured as events --- on \detone{} none of the 11 highest-multiplicity
events coincides with an HV spark --- so the HV log is the only reliable
discharge record, which is why this veto cannot be replaced by an event-level
cut.

\paragraph{Burst (event-level) veto.}
Localised discharges that are too small or too fast to move the integrated
mesh current above threshold still dump charge onto many pads at once. Since
real muon clusters have $n_\text{pad} \le 3$, any event with
$n_\text{pad} \ge 20$ is classified as a P2-internal discharge and vetoed.
The count of such bursts is logged per stage; the two vetoes together define
the ``spark-vetoed'' product variant, and every veto-capable stage is run
both with and without them so the size of the correction is always visible.

\paragraph{Per-pad micro-spark flagging (review only).}
A third, still finer population --- individual pads that micro-spark ---
survives both vetoes. These cannot be identified by rate alone (large outer
pads legitimately fire more often), so the discriminator is the
\emph{discharge amplitude signature}. For every pad with $\ge 20$ hits the
analysis computes the saturation rate and the fraction of hits above
\SI{1000}{ADC}, converts both to robust $z$-scores using the median/MAD
estimator $z = (x - \mathrm{med})/(1.4826\,\mathrm{MAD})$ over the live
pads, and flags a pad when \emph{both} are outliers
($z_\text{sat} > 4$ and $z_\text{hi} > 3$). A complementary rule flags
\emph{noise} pads first (median amplitude below \SI{150}{ADC} \emph{and}
firing-rate $z > 3$: high rate at near-pedestal amplitude), so the two
opposite-amplitude pathologies cannot contaminate each other's statistics.
Flagged pads are \emph{reviewed, not masked}: stage 08 renders what is
flagged and what the pad plane would look like without them, and the
decision to exclude them is taken explicitly, like the dead-connector drop.

\subsection{Efficiency measurement}
\label{sec:efficiency}

The efficiency chain (stage 06) proceeds strictly in the following order:

\begin{enumerate}[itemsep=2pt]
  \item \textbf{Channel table}: build the pad map for the run (validated
        \code{reverse} ordering, dead connectors dropped).
  \item \textbf{Rays}: load clean single-track M3 rays
        (Eq.~\eqref{eq:m3cut}) projected to the fitted plane $z$.
  \item \textbf{Centroids}: load per-event P2 centroids
        (Eq.~\eqref{eq:centroid}) and the set of events with \emph{any}
        mapped P2 hit.
  \item \textbf{Veto}: compute the spark/burst event-ID set and remove it
        from rays, centroids and the fired-event set alike.
  \item \textbf{Transform}: fit Eq.~\eqref{eq:procrustes} on the inner join
        of rays and centroids (fit fiducial $|x|,|y|<\SI{300}{mm}$), then map
        all centroids and all pad centres into the M3 frame.
  \item \textbf{Active area}: build a $k$-d tree of the transformed pad
        centres; a ray is \emph{in the active area} if its nearest
        transformed pad centre is within \SI{30}{mm}. This makes the
        denominator region fan-aware (it follows the true pad footprint,
        including the missing dead sectors) instead of a bounding box. The
        transformed pad footprint and the transformed insulation-mask pillar
        positions are persisted to CSV for the downstream sliding maps.
  \item \textbf{Per-ray flags}: for every in-active-area ray,
        \begin{itemize}
          \item $\code{within} \iff$ the event has a P2 centroid and
                $\lvert \mathbf{r}_\text{ray} - \mathbf{r}_\text{P2} \rvert
                 \le R$;
          \item $\code{has\_any} \iff$ the event contains any mapped P2 hit;
          \item $\code{resid} = \lvert \mathbf{r}_\text{ray} -
                \mathbf{r}_\text{P2}\rvert$ where defined.
        \end{itemize}
        A triggered ray whose event contains no P2 data at all is a genuine
        miss (\code{within} $=$ \code{has\_any} $=$ false) and \emph{stays in
        the denominator}: the trigger is external, and the DAQ writes an
        event whenever any detector (in particular the telescope itself)
        fired, so the absence of P2 hits is physical.
\end{enumerate}

The two ratios over active-area rays are the \emph{position-reconstruction
efficiency} (\code{within}, the physics number) and the \emph{response
ceiling} (\code{has\_any}); their difference, the ``fired but not
reconstructed'' fraction, isolates events where the detector responded but
the centroid landed farther than $R$ (edge charge sharing, delta rays,
mis-mapped pads). The match radius is not a free parameter tuned by eye: it
is chosen on the plateau of the efficiency-vs-$R$ curve of the validation
stage --- \SI{20}{mm} for \detone{}, \SI{40}{mm} for \dettwo{} --- and the
validation overlay separates the geometric containment cost of the cut,
$1-\exp\left(-R^2/2\sigma^2\right)$ for a Gaussian residual core of width
$\sigma$, from genuine detector loss. Products include the per-ray hit/miss
list (the input of stage 10), per-pad efficiency drawn on the true pad
tiles, the loss breakdown, the radial-residual distribution (core $+$ tail),
maps of the projections of unseen muons, and the efficiency in 30-minute
wall-clock bins with binomial errors --- the time series that exposed the
gain intermittency (Section~\ref{sec:results-det1}).

\subsection{Sliding-window efficiency maps}
\label{sec:sliding}

Hard 2D bins trade resolution against statistics once per map. Stage 10
instead evaluates, at every point of a fine regular grid, the mean of
\code{within} (and \code{has\_any}) over the rays inside a kernel:

\begin{itemize}[itemsep=2pt]
  \item \textbf{fixed kernel}: all rays within $\rho$ of the grid point
        ($\rho = \SI{5}{mm}$, minimum 10 rays), giving a smooth map whose
        resolution matches the $\sim$\SI{6}{mm} large insulation pillars ---
        overlapping kernels mean neighbouring pixels share rays, so
        structures are not aliased by bin edges;
  \item \textbf{adaptive $k$-NN kernel}: the kernel radius at each grid point
        is the distance to the $k$-th nearest ray, so every point averages
        exactly $k$ rays and carries the \emph{same} binomial uncertainty;
        the radius map is then itself the local-resolution map (points whose
        radius exceeds a cap are blanked).
\end{itemize}

Grid points whose nearest transformed pad lies farther than the footprint
radius are masked out, so the map is evaluated only on the physical detector
(using the footprint persisted by stage 06), and the transformed pillar
positions are overlaid so a pillar shadow is not misread as a gain defect.

\subsection{Efficiency and resolution vs.\ mesh voltage}
\label{sec:hvscan}

The HV scan (stage 11) reuses the stage-06 machinery with one structural
difference: the pad$\to$M3 transform is fitted \emph{once} on the pooled
matched events of \emph{all} HV points and then frozen, and with it the
transformed pad footprint and active area. The mounting geometry does not
change with voltage, so per-point re-fits would only inject alignment noise
--- particularly at low HV where few matched events exist --- and, worse,
would let the denominator region breathe across the scan. With the geometry
frozen, each sub-run contributes only its own rays and centroids:
efficiency (with binomial errors) and the core position resolution --- the
robust $\sigma$ (median/MAD) of the per-coordinate aligned residuals ---
are evaluated per HV point on an identical footing. Sub-runs without their
own \code{run\_config} borrow the long-run wiring; each sub-run's own
\code{hv\_monitor.csv} drives its spark veto.

\subsection{Pedestal QA}
\label{sec:pedestal}

Because the hit threshold is $5\sigma_\text{ped}$ \emph{per channel},
pedestal pathologies map directly onto rate pathologies, in opposite
directions: an anomalously low RMS lowers the threshold and lets the channel
over-trigger on baseline noise, while an inflated RMS (e.g.\ discharge
pickup during the pedestal run) silences real signal --- a ``cold'' pad that
looks dead without being broken. Stage 09 reads the pedestal means and RMS
values \emph{actually applied} by the processing (from the per-FEU
\code{hits\_root} files; the combined files drop them), maps both onto the
pad plane, flags robust outliers, and correlates pedestal RMS with the
per-pad firing rate. On \detone{} this correlation is negligible
($r \simeq 0.05$), which is the evidence that the observed hot pads are
genuine discharges (flat thresholds, anomalous amplitudes) rather than noisy
electronics --- the pedestal stage is what lets the pad-spark flagging
(Section~\ref{sec:sparks}) claim a physical cause.

\subsection{Structural validation of the efficiency}
\label{sec:validation}

Stage 12 exists to prove that the shipped efficiency is a property of the
detector, not of an analysis choice. It deliberately imports and calls the
\emph{same} functions as the production pipeline (the transform fit, the
loaders, the channel table), so the validator cannot diverge from what it
validates; every test has a predetermined pass condition, and a failure
points at a specific bias:

\begin{enumerate}[itemsep=2pt]
  \item \textbf{Knob plateau}: efficiency vs.\ match radius $R$ and vs.\
        active-area radius. An honest cut sits on a plateau; the $R$ curve is
        overlaid with the ideal Gaussian containment
        $1-\exp(-R^2/2\sigma^2)$, so the part of the loss that is pure cut
        geometry is visibly separated from detector loss.
  \item \textbf{$z$ plateau}: the whole chain (re-projection \emph{and}
        transform re-fit) is repeated across plane heights; the shipped $z$
        must sit on the efficiency plateau, otherwise efficiency is being
        left on the table by a wrong height.
  \item \textbf{Alignment variants}: free-scale fit vs.\ scale fixed to 1
        (pads are physical truth, so if forcing $s=1$ \emph{raises} the
        efficiency, the free fit was absorbing a bias) vs.\ a robust
        outlier-rejected re-fit; each variant reports its core $\sigma$ and
        efficiency.
  \item \textbf{Mapping vs.\ dead}: per-pad firing fraction against per-pad
        local efficiency. A real dead region is dark in \emph{both}; a
        mapping bug is dark in efficiency but bright in firing (the pads
        fire, but the reconstruction puts them elsewhere).
  \item \textbf{Estimator swap}: charge-weighted centroid vs.\
        nearest-fired-pad within $R$ --- bounds the cost of the centroiding
        choice itself.
  \item \textbf{Stability}: odd/even event split; the two halves must agree
        within statistics.
\end{enumerate}

For \dettwo{} six of seven tests pass; the single flag is an
\SI{8.2}{pt} sensitivity to the active-area radius (an edge-definition
systematic on a detector with several missing sectors), bounding the honest
efficiency between the shipped \SI{65.0}{\percent} and \SI{66.5}{\percent}
against the \SI{67.3}{\percent} any-pad ceiling.

% ========================================================================== %
\section{Results}
\label{sec:results}
% ========================================================================== %

\subsection{Detector 2 (\dettwo{})}
\label{sec:results-det2}

\dettwo{}, in its first run (10 h at mesh \SI{430}{V}), is the better-behaved
of the two detectors. The alignment is excellent: rotation
\SI{86.3}{\degree}, Procrustes scale $0.972$ (close to unity, as it must be
for a correct mapping), fit RMSE \SI{35.8}{mm}, fitted plane
$z=\SI{712}{mm}$. Table~\ref{tab:det2} summarises the long-run efficiency;
Fig.~\ref{fig:det2eff} shows the efficiency map and time evolution.

\begin{table}[htb]
  \centering
  \caption{\dettwo{} long-run results (mesh \SI{430}{V}, drift \SI{600}{V},
  spark-vetoed, connectors 1, 8--10 excluded).}
  \label{tab:det2}
  \small
  \begin{tabular}{lr}
    \toprule
    Clean M3 rays in active area & 59\,410 \\
    Efficiency (centroid within \SI{40}{mm}) & \SI{65.0}{\percent} \\
    Efficiency (any pad fired) & \SI{67.3}{\percent} \\
    Fired but not reconstructed & \SI{2.3}{\percent} \\
    Median $|r|$ residual & \SI{5.8}{mm} \\
    Core resolution $\sigma_x \times \sigma_y$ & $4.6 \times \SI{4.6}{mm}$ \\
    Validation & 6/7 pass; honest range 65.0--66.5\,\% \\
    HV sparks & 0.7/h, \SI{0.31}{\percent} dead time \\
    \bottomrule
  \end{tabular}
\end{table}

\begin{figure}[htb]
  \centering
  \begin{subfigure}[t]{.48\textwidth}
    \includegraphics[width=\textwidth]{det2/p2_det1_det2_long_run_mesh_scan_7-9-26/long_run_mesh_430V_drift_600V/06_efficiency/map_within_40mm_without_connectors_1_8_9_10_spark_vetoed.png}
    \caption{Binned efficiency map (centroid within \SI{40}{mm}).}
  \end{subfigure}\hfill
  \begin{subfigure}[t]{.48\textwidth}
    \includegraphics[width=\textwidth]{det2/p2_det1_det2_long_run_mesh_scan_7-9-26/long_run_mesh_430V_drift_600V/06_efficiency/efficiency_vs_time_without_connectors_1_8_9_10_spark_vetoed.png}
    \caption{Efficiency in 30-min bins over the \SI{10}{h} run.}
  \end{subfigure}
  \caption{\dettwo{} efficiency in the 9 July long run.}
  \label{fig:det2eff}
\end{figure}

Three observations qualify the headline number:

\begin{itemize}[itemsep=2pt]
  \item \textbf{Intermittency}: the 30-min binned efficiency spans
        0--\SI{75.5}{\percent}; the integrated \SI{65}{\percent} is a
        duty-cycle average. In its good periods the detector runs at
        $\sim$\SI{75}{\percent}, consistent with the HV-scan points below.
  \item \textbf{Sparking}: 7 mesh discharges in \SI{10}{h}
        (\SI{0.7}{/h}), dominated by one \SI{28}{s}, \SI{162}{\micro C}
        episode; the veto cost is only \SI{0.31}{\percent} of live time.
  \item \textbf{Resolution}: the core $\sigma\simeq\SI{4.7}{mm}$ is
        consistent with pure pad-pitch binning
        ($\text{pitch}/\sqrt{12}\simeq\SI{3.5}{mm}$ plus centroid sharing),
        i.e.\ the position response is healthy.
\end{itemize}

In the mesh scan (Fig.~\ref{fig:hvscans}, right) the efficiency rises
monotonically from \SI{65.1}{\percent} at \SI{395}{V} to
\SI{74.4}{\percent} at \SI{430}{V} with the resolution flat at
$\sigma\simeq\SI{4.7}{mm}$ --- the plateau has \emph{not} been reached at
\SI{430}{V}, and the scan should be extended upwards.

\subsection{Detector 1 (\detone{})}
\label{sec:results-det1}

\detone{} produced a solid first long run (30 June, mesh \SI{440}{V}):
integrated efficiency \SI{56.1}{\percent} (within \SI{20}{mm}),
any-pad \SI{64.8}{\percent}, median residual \SI{11}{mm} over 75\,035 active
rays. Already there, however, the efficiency-vs-time trace was strongly
intermittent (0--\SI{74}{\percent} per 30-min bin), so the integrated number
is a duty-cycle average of a detector that alternates between a
$\sim$\SI{74}{\percent} ``on'' state and a silent state. Across subsequent
runs at the same working point the duty cycle collapsed
(Table~\ref{tab:det1runs}); in the 9~July combined run (\SI{10}{V} lower mesh)
the detector was essentially inefficient. The HV-spark rate on \detone{} is
negligible (0.1/h), so discharges do not explain the drop-outs.

\begin{table}[htb]
  \centering
  \caption{\detone{} integrated efficiency at the nominal working point across
  the campaign (spark-vetoed). The range column is the 30-min binned
  \code{within} efficiency, exposing the intermittency.}
  \label{tab:det1runs}
  \small
  \begin{tabular}{llrrl}
    \toprule
    Run & Mesh & Within-$R$ & Any-pad & 30-min range \\
    \midrule
    6-30-26 & \SI{440}{V} & \SI{56.1}{\percent} & \SI{64.8}{\percent} & 0--\SI{74.1}{\percent} \\
    7-4-26  & \SI{440}{V} & \SI{3.5}{\percent}  & \SI{8.3}{\percent}  & 0.1--\SI{20.1}{\percent} \\
    7-7-26  & \SI{440}{V} & \SI{11.5}{\percent} & \SI{15.8}{\percent} & 3.4--\SI{62.1}{\percent} \\
    7-9-26  & \SI{430}{V} & \SI{0.4}{\percent}  & \SI{3.0}{\percent}  & 0--\SI{0.8}{\percent} \\
    \bottomrule
  \end{tabular}
\end{table}

Two further \detone{}-specific findings from the good run:

\begin{itemize}[itemsep=2pt]
  \item a real \emph{diagonal near-zero-efficiency band} crosses the fan
        (reference $x \sim$ 120--\SI{160}{mm}), dark in both firing and
        efficiency (hence a genuine dead/low-gain region or mapping seam, not
        a reconstruction artefact); it has not yet been traced to specific
        pads or connectors;
  \item a population of pads clustered in connectors 6/7 carries a per-pad
        micro-discharge signature (saturation + high-amplitude outliers);
        pedestal QA confirms these are real discharges, not noisy channels.
\end{itemize}

The 2 July HV scan (Fig.~\ref{fig:hvscans}, left), taken while the detector
was behaving, shows the expected S-curve: efficiency rising from
\SI{16.5}{\percent} at \SI{345}{V} to \SI{68.4}{\percent} at \SI{420}{V},
still without a plateau, with resolution roughly constant at
$\sigma\sim\SI{10}{mm}$.

\begin{figure}[htb]
  \centering
  \begin{subfigure}[t]{.48\textwidth}
    \includegraphics[width=\textwidth]{det1/p2_det1_mesh_hv_scan_7-2-26/hv_scan/11_hv_scan_efficiency/efficiency_vs_hv_without_connectors_1_10_spark_vetoed.png}
    \caption{\detone{}, 2 July scan (mesh 345--\SI{420}{V}).}
  \end{subfigure}\hfill
  \begin{subfigure}[t]{.48\textwidth}
    \includegraphics[width=\textwidth]{det2/p2_det1_det2_long_run_mesh_scan_7-9-26/hv_scan/11_hv_scan_efficiency/efficiency_vs_hv_without_connectors_1_8_9_10_spark_vetoed.png}
    \caption{\dettwo{}, 9 July scan (mesh 395--\SI{430}{V}).}
  \end{subfigure}
  \caption{Efficiency vs.\ mesh HV. Neither detector has reached its plateau
  at the highest scanned voltage.}
  \label{fig:hvscans}
\end{figure}

% ========================================================================== %
\section{Conclusions}
\label{sec:conclusions}
% ========================================================================== %

\begin{itemize}[itemsep=3pt]
  \item The full measurement chain --- automated DAQ, on-the-fly processing,
        M3 reference tracking and the staged QA pipeline --- is operational
        and has characterised two P2 pad detectors end to end, including a
        structural validation of the efficiency methodology itself.
  \item The Gerber-derived pad mapping, combined with the validated
        \code{reverse} within-connector channel ordering, places hits on the
        physically correct pads for \emph{both} detectors (track--centroid
        correlations up to $r_x = 0.93$, transform scale
        $\simeq 1$); the ordering is a fixed property of the DREAM readout
        and can be trusted for future P2 detectors.
  \item \dettwo{} performs well: \SI{65}{\percent} integrated
        (74--\SI{75}{\percent} in good periods and at the best HV points),
        core resolution $\sigma\simeq\SI{4.7}{mm}$ (pad-pitch limited), a
        validated efficiency estimate, and a modest, cleanly vetoed sparking
        rate.
  \item \detone{} is limited not by its intrinsic response (it reaches
        $\sim$\SI{74}{\percent} when ``on'') but by \emph{intermittent gain
        drop-outs} whose duty cycle degraded run after run, down to
        near-total inefficiency by 9 July. HV sparks are ruled out as the
        cause; the diagonal dead band and the connector-6/7 micro-discharge
        population are additional, distinct defects.
  \item Neither HV scan reached the efficiency plateau
        (\detone{} scanned to \SI{420}{V}, \dettwo{} to \SI{430}{V}); current
        working points sit on the rising edge of the S-curve.
\end{itemize}

% ========================================================================== %
\section{Roadmap and future improvements}
\label{sec:roadmap}
% ========================================================================== %

\subsection{Detector and hardware}
\begin{enumerate}[itemsep=2pt]
  \item \textbf{Diagnose the \detone{} intermittency} (top priority):
        correlate the efficiency-vs-time drop-outs with the full HV log
        (voltage sag, trips, current baseline), gas flow and environmental
        records; test whether HV retraining or a gas purge restores the
        duty cycle; if not, inspect the detector (drift contact, mesh
        polarisation).
  \item \textbf{Recover the dead channels}: re-connect connector~1 on both
        detectors and debug the FEU~8 crash that costs \dettwo{} connectors
        8--10 (three of ten sectors).
  \item \textbf{Extend the HV scans upwards} (\dettwo{} beyond \SI{430}{V},
        \detone{} beyond \SI{440}{V} once stable) with the spark rate
        monitored as the limiting variable, to locate the plateau and the
        efficiency--discharge trade-off; longer sub-runs per point would cut
        the $\sim$\SI{1}{\percent} statistical error per point.
  \item \textbf{Trace the \detone{} diagonal dead band} to hardware: overlay
        the band on the connector/sector geometry and the insulation mask,
        and probe the corresponding pads/connectors on the bench.
  \item \textbf{Systematic HV training} of new detectors with the existing
        \code{hv\_trainer} tooling before physics runs.
\end{enumerate}

\subsection{Analysis and pipeline}
\begin{enumerate}[itemsep=2pt]
  \item \textbf{Report the ``on-state'' efficiency}: separate the duty cycle
        from the intrinsic efficiency by gating on the good 30-min bins (or
        an HMM-style on/off classification), so detector physics and
        stability are quoted as two distinct numbers.
  \item \textbf{Close the last validation flag}: make the \dettwo{} active-area
        (edge) definition robust, e.g.\ by an explicit fiducial contour with
        an edge-distance systematic, removing the \SI{8.2}{pt} sensitivity.
  \item \textbf{Pillar-level efficiency}: exploit the \SI{5}{mm} sliding maps
        and the insulation-mask overlay to quantify the local efficiency
        loss at the five large pillars, and grow statistics until the small
        pillar grid can be tested.
  \item \textbf{Amplitude and gain analysis}: extend the QA with per-pad gain
        (Landau MPV) maps and gain-vs-HV curves from the scan data, giving a
        gain-based (not just efficiency-based) working-point choice.
  \item \textbf{Two-detector combined analysis}: with both planes aligned to
        M3, measure \dettwo{} efficiency using \detone{}-confirmed tracks
        (and vice versa) as a cross-check of the telescope-based numbers,
        and study inter-detector timing.
  \item \textbf{Automation}: fold stages 10--12 into the on-the-fly QA so a
        validated efficiency and its time trace are available during the run,
        allowing bad periods (like the \detone{} drop-outs) to be caught and
        acted on live.
\end{enumerate}

\vspace{2ex}
\noindent\textit{All numbers in this note are reproducible from the run
registry in \code{p2\_qa\_config.py}: run
\code{./run\_p2\_pipeline.sh $\langle$run\_key$\rangle$} followed by
\code{python3 build\_final\_pdf.py $\langle$run\_key$\rangle$}; the HV scans
use stage 11 and \code{build\_hv\_scan\_pdf.py}.}

\end{document}
